El experimento demuestra que el overclocking factor debe seleccionarse considerando la frecuencia interna de los módulos para lograr resultados coherentes en FIL. Para el contador evaluado, el valor de \texttt{OF} que mejor reproduce un incremento de 100 en 1 segundo es aquel que coincide con la frecuencia de conteo del canal observado. En aplicaciones con acumulaciones internas no limitadas, un \texttt{OF} mayor incrementa los resultados observados y puede llevar a interpretaciones erróneas si no se justifica correctamente el escalado temporal.

Los resultados confirman que el factor de overclocking controla cuántas actualizaciones internas del FPGA quedan encapsuladas en cada paso de Simulink. Si el diseño no tiene mecanismos de acumulación que limiten su crecimiento, un \texttt{OF} alto produce valores observados mucho mayores a los esperados para una escala temporal fija en Simulink. Esto es evidente en los contadores de 1 y 10 ciclos cuando se usa \texttt{OF = 17}. Asimismo, se observa que el muestreo de salida introduce una cuantización temporal: cuando la salida se actualiza en el FPGA con una frecuencia distinta a la de muestreo de Simulink, puede aparecer imprecisión aparente. Por ejemplo, con \texttt{OF = 10}, el contador de 17 ciclos se actualiza a un ritmo más lento que el muestreo, lo que genera saltos visibles y un valor final inferior a 100.
